% AI 使用说明。正式提交前请根据最终保留的调用记录核对工具和模型信息。
\section{人工智能工具使用说明}
\label{app:ai-use}

本研究在资料整理、字段与来源核对、候选程序与绘图框架、数值一致性检查以及
\LaTeX{} 结构整理等环节使用了由 OpenAI 提供的 Codex desktop 作为辅助工具，
不替代队员的独立思考和核心工作。

\begin{description}
  \item[资料与程序辅助：]
  根据既定数据契约提供字段核对、代码框架、调试和绘图建议。数据清洗规则、
  程序逻辑、参数设置和运行结果由队员审定并复现。

  \item[分析与复核辅助：]
  对已生成结果提供候选重算、敏感性检查和一致性检查建议。模型选择、阈值设定、
  证据等级判断以及结果解释由队员讨论、审定并负责。

  \item[写作与排版辅助：]
  提供段落结构、图注和排版方面的候选建议。最终文字由队员理解后用自己的语言
  重写，并逐项核对数据、公式和引用。
\end{description}

研究问题、指标定义、关键假设、模型结构、参数选择、实验设计和最终结论均由队员
提出、讨论和审定。拟纳入最终论文主张的代码、公式、数值和文字均经过队员理解、
修改、脚本重跑或交叉核对；未通过复核、证据不足或仍处于探索性/待评审状态的内容，
不作为确定性结论。正文中的模型、公式、数据和参考文献均以可核验的推导、实验记录
和正式文献为依据，人工智能工具输出不作为独立证据来源。使用过程中未向公共 AI 服务
提交未公开题目、原始数据或私密对话。

\noindent 详细的提示词、调用时间、输入摘要、输出引用、运行号和人工复核状态记录于
项目的 \texttt{governance/ai-use-log.csv} 及相关项目材料中。
